%
% teil4.tex -- Anwendung: Lichtbeugung am Kreisspalt
%
% (c) 2026 Gian Cavegn, OST Ostschweizer Fachhochschule
%
% !TEX root = ../../buch.tex
% !TEX encoding = UTF-8
%
\section{Anwendung: Lichtbeugung am Kreisspalt
\label{bessel:section:beugung}}
\kopfrechts{Lichtbeugung}
% Funktion dieses Abschnitts:
% Erste Anwendung der entwickelten Theorie. Zeigt, wie die
% rotationssymmetrische Geometrie der kreisförmigen Blende direkt
% auf eine Bessel-Funktion führt -- das Airy-Pattern.
% Kürzer halten als die mathematischen Abschnitte (1-2 Seiten).

Als erste Anwendung betrachten wir die Beugung von Licht an einer kreisförmigen Öffnung.
Dieses Problem ist von grundlegender Bedeutung in der Optik, da es das Auflösungsvermögen jedes optischen Instruments mit kreisförmiger Apertur (Teleskop, Mikroskop, Auge) begrenzt.

\subsection{Vom Einzelspalt zum Kreis
\label{bessel:subsection:vomspalt}}
% Inhalt:
% - Erinnerung: bei einem rechteckigen Spalt ergibt die Fraunhofer-
%   Beugung eine sinc-Funktion.
% - Bei einem kreisförmigen Spalt wird der Ansatz analog, nur in
%   Polarkoordinaten in der Blendenebene.
% - Daraus entsteht das 2D-Integral der Beugungsformel.

Bei einem rechteckigen Spalt der Breite $b$ ergibt die Fraunhofer-Beugung eine Intensitätsverteilung, die durch die $\operatorname{sinc}$-Funktion beschrieben wird.
Der Übergang von einer rechteckigen zu einer kreisförmigen Öffnung erfordert es, das Beugungsintegral in Polarkoordinaten der Blendenebene zu formulieren.

\subsection{Das Beugungsintegral
\label{bessel:subsection:beugungsintegral}}
% Inhalt:
% - Fraunhofer-Beugungsintegral in 2D hinschreiben.
% - Übergang in Polarkoordinaten r, phi in der Blendenebene.
% - Die phi-Integration liefert direkt eine Bessel-Funktion:
%   \int_0^{2\pi} e^{i k r sin(theta) cos(phi)} d\phi = 2\pi J_0(k r sin theta)
% - Anschliessende r-Integration über [0, R] mit der Identität
%   \int x J_0(x) dx = x J_1(x).

Die Fraunhofer-Beugungsamplitude in der Brennebene einer Linse hinter einer Apertur lautet
\begin{equation}
U(\theta)
=
C \int_{\text{Apertur}} e^{-i\vec k \cdot \vec r}\, d^2 r,
\label{bessel:equation:fraunhofer}
\end{equation}
mit einer geometrieabhängigen Konstanten $C$.
In Polarkoordinaten $\vec r = (r\cos\varphi, r\sin\varphi)$ und für eine kreisförmige Apertur vom Radius $R$ wird daraus
\begin{equation}
U(\theta)
=
C \int_0^R \int_0^{2\pi} e^{-i k r \sin\theta\, \cos\varphi}\, r\, d\varphi\, dr.
\end{equation}
Das innere $\varphi$-Integral ergibt mit der Integraldarstellung der Bessel-Funktion
\begin{equation}
\int_0^{2\pi} e^{-i z \cos\varphi}\, d\varphi
=
2\pi\, J_0(z)
\label{bessel:equation:integral-J0}
\end{equation}
% [TODO: Diese Integraldarstellung von J_0 müsste in Abschnitt~3 oder
% hier kurz hergeleitet werden]
den Wert $2\pi J_0(k r \sin\theta)$.
Mit der aus~\eqref{bessel:equation:rek2} und~\eqref{bessel:equation:J0-ableitung} ableitbaren Identität $\int_0^R x J_0(\alpha x)\, dx = (R/\alpha) J_1(\alpha R)$ folgt schliesslich
\begin{equation}
U(\theta)
\propto
\frac{J_1(k R \sin\theta)}{k R \sin\theta}.
\label{bessel:equation:amplitude}
\end{equation}

\subsection{Das Airy-Pattern
\label{bessel:subsection:airy}}
% Inhalt:
% - Intensität ist |U|^2.
% - Das Resultat heisst Airy-Funktion bzw. Airy-Pattern.
% - Erstes Minimum bei k R sin(theta) = j_{1,1} = 3.8317.
% - Daraus ergibt sich das Rayleigh-Kriterium fuer das
%   Auflösungsvermögen: sin(theta) = 1.22 lambda / D, mit D = 2R.

Die beobachtete Intensität ist proportional zum Betragsquadrat der Amplitude:
\begin{equation}
I(\theta)
=
I_0 \left[\frac{2 J_1(k R \sin\theta)}{k R \sin\theta}\right]^2.
\label{bessel:equation:airy}
\end{equation}
Diese Funktion heisst \emph{Airy-Funktion}; das zugehörige Beugungsmuster -- ein helles Zentralmaximum, umgeben von konzentrischen Ringen -- ist als \emph{Airy-Scheibchen} bekannt.

% [TODO: Bild des Airy-Patterns einfügen]

Das erste Minimum tritt dort auf, wo das Argument der Bessel-Funktion gleich der ersten Nullstelle $j_{1,1} \approx 3{,}8317$ ist, also bei
\begin{equation}
\sin\theta_{\min}
=
\frac{j_{1,1}}{kR}
=
\frac{j_{1,1}\, \lambda}{2\pi R}
\approx
1{,}22 \cdot \frac{\lambda}{D},
\label{bessel:equation:rayleigh}
\end{equation}
wobei $D = 2R$ der Durchmesser der Apertur und $\lambda$ die Wellenlänge ist.
Diese Beziehung ist das berühmte \emph{Rayleigh-Kriterium} für das Auflösungsvermögen optischer Instrumente.
% [TODO: Kurzer Kommentar zur praktischen Bedeutung: Teleskope, Mikroskope, Auge]

